\documentclass[11pt]{article}

\usepackage[margin=1in]{geometry}
\usepackage{amsmath, amssymb}
\usepackage{booktabs}
\usepackage{graphicx}
\usepackage[colorlinks=true, linkcolor=blue, urlcolor=blue]{hyperref}

\newcommand{\istd}{\texttt{intrinsic\_std}}
\newcommand{\irms}{\texttt{intrinsic\_rms}}
\newcommand{\imad}{\texttt{intrinsic\_mad}}
\newcommand{\moff}{\texttt{mean\_median\_offset}}
\newcommand{\gp}{\gamma_{\mathrm{P}}}
\newcommand{\sigobs}{\sigma_{\mathrm{obs}}}
\newcommand{\sigphot}{\sigma_{\mathrm{phot}}}
\newcommand{\sigint}{\sigma_{\mathrm{int}}}

\title{Variability Summary Statistics for the TESS Cluster Table:\\
       Noise Correction, Normalisation, and Sector Stitching}
\author{TESS Cluster Age ML --- \texttt{GenerateTable\_CGW}}
\date{2026 August 19}

\begin{document}
\maketitle

\begin{abstract}
This note documents the six per-row variability statistics computed by
\texttt{table\_pipeline/summary\_stats.py} and analyses each along three axes:
whether it admits an exact photometric-noise correction, whether it is invariant
under the pipeline's median flux normalisation, and how it behaves when several
TESS sectors are stitched into a single row. The organising result is that a
statistic is analytically noise-correctable if and only if it is built from
\emph{cumulants}, which add under convolution; statistics built from
\emph{quantiles} (the median, the MAD) are not correctable in closed form. This
cleanly separates the three amplitude statistics and the moment skewness
$g_1$, which are correctable, from the MAD-based amplitude and the Pearson
nonparametric skew $\gp$, which are not. A second, independent axis --- stability
under sector mixing --- ranks the statistics in the opposite order, so no single
statistic dominates.
\end{abstract}

\tableofcontents

% ============================================================================
\section{Notation and pipeline context}
% ============================================================================

\subsection{Construction of a table row}

Each row of the cluster table corresponds to one cluster observed in one
\emph{combination} of TESS sectors. The light curve for a row is assembled in
\texttt{lc\_lsp.py:load\_and\_process\_sectors} by the following sequence, applied
independently to each sector $k$ in the combination:

\begin{enumerate}
  \item drop non-finite samples;
  \item block-average (\texttt{resample\_lc}) to a common $30$\,min cadence,
        performed on \emph{raw} flux, discarding any block spanning a time gap;
  \item MAD-based sigma clipping (\texttt{mad\_clip\_lc}) at $n_\sigma = 4$,
        applied \emph{symmetrically} about the sector median;
  \item median normalisation (\texttt{normalize\_flux}), see below;
\end{enumerate}

after which the normalised sectors are concatenated and sorted in time.

\subsection{The median normalisation}

For a single sector with raw flux $f_i$ and uncertainty $\varsigma_i$,
\texttt{normalize\_flux} applies
%
\begin{equation}
  x_i = \frac{f_i}{\mathrm{med}(f)}, \qquad
  s_i = \frac{\varsigma_i}{\mathrm{med}(f)} .
  \label{eq:norm}
\end{equation}
%
This is a purely \emph{multiplicative} rescaling: it is not a mean subtraction.
Two consequences are used repeatedly below.

\paragraph{(i) The sector median is exactly unity.}
By construction $\mathrm{med}(x) = \mathrm{med}(f)/\mathrm{med}(f) = 1$.

\paragraph{(ii) The \emph{pooled} median is also exactly unity.}
Each sector contributes exactly half of its own points below its own median.
Summing over $K$ sectors with $n_k$ points each, the concatenated array has
$\sum_k n_k/2 = N/2$ points below $1$, so the pooled median is $1$ regardless of
how unequal the sector lengths are (up to the usual $\pm 1$ discreteness for odd
$n_k$). This was verified numerically on the full table: over $1250$ rows,
$\max|\mathrm{med}(x) - 1| = 2.2\times10^{-16}$, and separately for
$n_{\rm sec} = 1$, $2$, and $3$.

Because both the flux and its uncertainty are divided by the same constant, all
statistics below are in dimensionless \emph{relative flux} units, which is what
makes them comparable between clusters of different apparent brightness.

\subsection{Basic quantities}

For a row with $N$ samples $\{x_i\}$ and uncertainties $\{s_i\}$, write
%
\begin{align}
  \bar{x}       &= \frac{1}{N}\sum_{i=1}^{N} x_i,
  &
  \sigphot^{2}  &= \frac{1}{N}\sum_{i=1}^{N} s_i^{2},
  \label{eq:basic1}\\[4pt]
  \sigobs^{2}   &= \frac{1}{N}\sum_{i=1}^{N}\left(x_i - \bar{x}\right)^{2},
  &
  \mu_3         &= \frac{1}{N}\sum_{i=1}^{N}\left(x_i - \bar{x}\right)^{3}.
  \label{eq:basic2}
\end{align}
%
Note that $\sigphot^2$ is the \emph{mean of the squared} error bars, not the
square of the mean error bar, and that all central moments use $\mathrm{ddof}=0$.

% ============================================================================
\section{The noise model and the governing principle}
% ============================================================================

\subsection{Model}

Write the observed flux as the sum of an astrophysical signal and a measurement
error,
%
\begin{equation}
  x = s + n,
  \label{eq:model}
\end{equation}
%
with $n$ independent of $s$, zero-mean, and \emph{symmetric}, with variance
$\mathrm{Var}(n) = \sigphot^2$. Symmetry is the operative assumption: TESS
photometric errors are approximately Gaussian, and Gaussians are symmetric.

\subsection{Cumulants add; quantiles do not}

Equation~\eqref{eq:model} is a convolution of probability densities. Under
convolution of independent variables, \emph{cumulants} $\kappa_j$ are additive:
%
\begin{equation}
  \kappa_j(s+n) = \kappa_j(s) + \kappa_j(n).
  \label{eq:cumadd}
\end{equation}
%
The first three cumulants are the mean, the variance, and the third \emph{central
moment} $\mu_3$ (for $j\le3$ cumulant and central moment coincide). Applying
\eqref{eq:cumadd} to each:
%
\begin{align}
  \bar{x}       &= \bar{s} + \underbrace{E[n]}_{=\,0} = \bar{s}
                 && \text{(immune)},
  \label{eq:cum1}\\
  \sigobs^{2}   &= \sigma_s^{2} + \sigphot^{2}
                 && \text{(inflated, exactly correctable)},
  \label{eq:cum2}\\
  \mu_3(x)      &= \mu_3(s) + \underbrace{\mu_3(n)}_{=\,0} = \mu_3(s)
                 && \text{(immune, since $n$ is symmetric)}.
  \label{eq:cum3}
\end{align}
%
Equation~\eqref{eq:cum3} is the central result of this note: \textbf{symmetric
noise contributes exactly zero third moment}, so the numerator of any
moment-based skewness requires no noise subtraction at all. Only the
normalising $\sigma^3$ is contaminated, and \eqref{eq:cum2} tells us exactly how.

By contrast, \emph{quantiles} are not cumulants and do not obey
\eqref{eq:cumadd}. In particular
%
\begin{equation}
  \mathrm{med}(s+n) \;\neq\; \mathrm{med}(s)
  \label{eq:medshift}
\end{equation}
%
in general. Convolving a skewed density with symmetric noise fills in the sparse
side of the distribution and \emph{pulls the median toward the mean} by an amount
that depends on the unknown shape of $p(s)$. There is no closed-form inverse.
The same objection applies to the median absolute deviation.

\paragraph{The unavoidable trade-off.} The property that makes the median robust
to outliers --- that it counts points rather than weighting them by magnitude ---
is precisely the property that breaks additivity under convolution. Robustness to
outliers and analytic noise-correctability are therefore in direct tension: among
the statistics below, no quantile-based estimator is exactly correctable, and no
exactly correctable estimator is outlier-robust.

% ============================================================================
\section{The statistics}
% ============================================================================

Throughout, $\sigint \equiv \istd$ denotes the noise-corrected scale defined in
\S\ref{sec:std}. Every ``intrinsic'' quantity in the pipeline applies the floor
$\max(0,\cdot)$ before taking a square root; the consequences are discussed in
\S\ref{sec:floor}.

% ----------------------------------------------------------------------------
\subsection{Noise-corrected standard deviation, \istd}
\label{sec:std}
% ----------------------------------------------------------------------------

\paragraph{Base form.} $\sigobs$ as defined in \eqref{eq:basic2}: the scatter
about the row's own mean.

\paragraph{Noise-corrected form.} Directly from \eqref{eq:cum2},
%
\begin{equation}
  \boxed{\;\istd \;=\; \sigint \;=\; \sqrt{\max\!\left(0,\;
         \sigobs^{2} - \sigphot^{2}\right)}\;}
  \label{eq:istd}
\end{equation}
%
This correction is \textbf{exact}. The variance is a cumulant, the subtraction is
the identity \eqref{eq:cum2} rearranged, and no distributional assumption beyond
independence is required.

% ----------------------------------------------------------------------------
\subsection{Noise-corrected RMS, \irms}
\label{sec:rms}
% ----------------------------------------------------------------------------

\paragraph{Base form.} The scatter about the \emph{absolute} baseline $1$,
%
\begin{equation}
  \mathrm{rms}_{\mathrm{obs}}^{2} = \frac{1}{N}\sum_{i}\left(x_i - 1\right)^{2}.
  \label{eq:rmsobs}
\end{equation}

\paragraph{Relation to \istd.} Writing $x_i - 1 = (x_i - \bar{x}) + (\bar{x}-1)$
and expanding, the cross term vanishes because $\sum_i (x_i - \bar{x}) = 0$:
%
\begin{equation}
  \mathrm{rms}_{\mathrm{obs}}^{2}
    = \sigobs^{2} + \left(\bar{x} - 1\right)^{2}
    = \sigobs^{2} + \Delta^{2},
  \label{eq:rmsidentity}
\end{equation}
%
where $\Delta \equiv \bar{x} - \mathrm{med}(x)$ is the offset of
\S\ref{sec:offset}, the substitution $1 \to \mathrm{med}(x)$ being exact by the
normalisation. Subtracting $\sigphot^2$ from both sides,
%
\begin{equation}
  \boxed{\;\irms^{2} \;=\; \istd^{2} + \Delta^{2}\;}
  \label{eq:rmsstd}
\end{equation}
%
valid wherever neither quantity is truncated by its $\max(0,\cdot)$ floor. This
was verified on real data: over the $1216$ of $1250$ rows where neither is
clamped, \eqref{eq:rmsstd} holds to a maximum absolute residual of
$2.9\times10^{-19}$.

\paragraph{Interpretation.} \irms{} carries \emph{no amplitude information beyond}
\istd. The two are degenerate for any symmetric light curve, and their difference
is purely the squared asymmetry. This makes \irms{} a hybrid object: the
$\sigobs^2$ part of \eqref{eq:rmsidentity} is a cumulant and is exactly
corrected, but the $\Delta^2$ part is quantile-contaminated and is \emph{not}.
Its noise correction is therefore exact only in the symmetric limit
$\Delta \to 0$.

% ----------------------------------------------------------------------------
\subsection{Noise-corrected MAD, \imad}
\label{sec:mad}
% ----------------------------------------------------------------------------

\paragraph{Base form.} The median absolute deviation rescaled to a
Gaussian-equivalent $\sigma$,
%
\begin{equation}
  \mathrm{MAD} = \mathrm{med}\left(\left|x_i - \mathrm{med}(x)\right|\right),
  \qquad
  \sigma_{\mathrm{MAD}} = 1.4826\,\mathrm{MAD},
  \label{eq:mad}
\end{equation}
%
where $1.4826 = 1/\Phi^{-1}(3/4)$ is the consistency factor making
$\sigma_{\mathrm{MAD}}$ an unbiased estimator of $\sigma$ \emph{for Gaussian
data}.

\paragraph{Nominal noise-corrected form.}
%
\begin{equation}
  \imad = \sqrt{\max\!\left(0,\;\sigma_{\mathrm{MAD}}^{2}
                            - \sigphot^{2}\right)} .
  \label{eq:imad}
\end{equation}

\paragraph{Caveat --- this correction is an approximation, not an identity.}
Equation~\eqref{eq:imad} presumes that $\sigma_{\mathrm{MAD}}^2$ is additive
under convolution. It is not: the MAD is a quantile statistic and violates
\eqref{eq:cumadd}. Additivity holds exactly only when \emph{both} $s$ and $n$
are Gaussian, in which case $\sigma_{\mathrm{MAD}} = \sigma$ identically and one
may as well have used \eqref{eq:istd}. For genuinely non-Gaussian signal --- spot
modulation, flares, precisely the regime in which a robust estimator is wanted ---
\eqref{eq:imad} carries a shape-dependent bias. \imad{} should be read as a
robust \emph{diagnostic} rather than as an exactly de-biased amplitude, and
disagreement between \imad{} and \istd{} is itself informative about
non-Gaussianity.

% ----------------------------------------------------------------------------
\subsection{Mean--median offset, \moff}
\label{sec:offset}
% ----------------------------------------------------------------------------

\paragraph{Base form.}
%
\begin{equation}
  \boxed{\;\Delta \;\equiv\; \moff \;=\; \bar{x} - \mathrm{med}(x)\;}
  \label{eq:offset}
\end{equation}

\paragraph{Noise-corrected form: none exists.} By \eqref{eq:cum1} the mean is
immune, but by \eqref{eq:medshift} the median is not: as noise increases, the
observed median migrates toward the observed mean and $\Delta$ collapses toward
zero. Since the shift depends on the unknown shape of $p(s)$, no closed-form
correction is available. Note that this is a bias in the \emph{numerator}, so it
cannot be repaired by any change of denominator.

\paragraph{Dimensions.} $\Delta$ has units of relative flux; it is an
\emph{amplitude}, not a shape parameter, and belongs with \istd{}/\irms{}/\imad{}
rather than with $\gp$ and $g_1$.

% ----------------------------------------------------------------------------
\subsection{Pearson nonparametric skew, $\gp$}
\label{sec:gammap}
% ----------------------------------------------------------------------------

\paragraph{Base form.}
%
\begin{equation}
  \boxed{\;\gp \;=\; \frac{\bar{x} - \mathrm{med}(x)}{\sigobs}
                 \;=\; \frac{\Delta}{\sigobs}\;}
  \label{eq:gammap}
\end{equation}
%
defined as \texttt{NaN} where $\sigobs = 0$. Negative $\gp$ indicates a
dip-dominated (spot-like) flux distribution; positive $\gp$ a flare-dominated
one.

\paragraph{Mechanism.} $\gp$ is a \emph{location-contrast} statistic. It infers
asymmetry from the disagreement between two location estimators: mean and median
coincide for a symmetric density and separate when a tail drags the mean.

\paragraph{Partially corrected form.} One may substitute the noise-corrected
scale in the denominator,
%
\begin{equation}
  \gp^{\mathrm{(int)}} = \frac{\Delta}{\sigint}
                       = \gp \cdot \frac{\sigobs}{\sigint},
  \label{eq:gammapint}
\end{equation}
%
but this repairs only the denominator. Because $\Delta$ itself is biased, the
result is \textbf{not} an unbiased estimate of the signal's $\gp$. An injection
test (Table~\ref{tab:injection}) confirms this: applying \eqref{eq:gammapint} to
data with a known true $\gp = -0.640$ recovers $-0.545$ at $\mathrm{SNR}=10$ and
only $-0.163$ at $\mathrm{SNR}=1$ --- essentially identical to the uncorrected
value, because the denominator was never the problem.

% ----------------------------------------------------------------------------
\subsection{Moment skewness, $g_1$}
\label{sec:g1}
% ----------------------------------------------------------------------------

\paragraph{Base form.}
%
\begin{equation}
  g_1 = \frac{\mu_3}{\sigobs^{3}}
      = \frac{N^{-1}\sum_i \left(x_i-\bar{x}\right)^{3}}{\sigobs^{3}} .
  \label{eq:g1}
\end{equation}

\paragraph{Mechanism.} $g_1$ is a \emph{moment} statistic. Cubing the deviations
preserves sign, so points far above the mean contribute positively and points far
below negatively; a symmetric density cancels term by term. Unlike $\gp$, it
never references the median.

\paragraph{Noise-corrected form.} By \eqref{eq:cum3} the numerator needs no
correction; by \eqref{eq:cum2} the denominator is corrected by replacing
$\sigobs$ with $\sigint$. The inflation factors cancel and the result collapses
to a single expression:
%
\begin{equation}
  \boxed{\;g_1^{\mathrm{(int)}}
     \;=\; g_1 \left(\frac{\sigobs}{\sigint}\right)^{3}
     \;=\; \frac{\mu_3}{\istd^{3}}\;}
  \label{eq:g1int}
\end{equation}
%
That is: \textbf{the noise-corrected skewness is the third central moment divided
by the cube of the already-computed} \istd. Structurally this is exactly parallel
to \eqref{eq:istd} --- the same numerator, normalised by the noise-corrected scale
instead of the observed one. The correction is \emph{exact} under the symmetric
noise assumption.

\paragraph{Numerical verification.} Table~\ref{tab:injection} reports an
injection test: an asymmetric spot-like signal of known skewness
($g_1^{\rm true} = -1.3087$, $\mu_3^{\rm true} = -1.3087$ at unit variance,
$\Delta^{\rm true} = -0.640$) with white noise added at decreasing SNR,
$N = 2000$ samples, $200$ realisations per SNR.

\begin{table}[htbp]
\centering
\caption{Injection test, mean recovered value over $200$ realisations. The
$\mu_3$ column is flat across a tenfold change in noise, confirming
\eqref{eq:cum3}; $g_1^{\mathrm{(int)}}$ recovers truth to four significant
figures at every SNR. The $\Delta$ column collapses by a factor of four,
confirming \eqref{eq:medshift}, and $\gp^{\mathrm{(int)}}$ inherits that
collapse.}
\label{tab:injection}
\begin{tabular}{r rr rr}
\toprule
 & \multicolumn{2}{c}{moment route} & \multicolumn{2}{c}{quantile route} \\
\cmidrule(lr){2-3}\cmidrule(lr){4-5}
SNR & $\mu_3$ & $g_1^{\mathrm{(int)}}$ & $\Delta$ & $\gp^{\mathrm{(int)}}$ \\
\midrule
truth   & $-1.3087$ & $-1.3087$ & $-0.6399$ & $-0.6399$ \\
\midrule
$10.0$  & $-1.3086$ & $-1.3087$ & $-0.5452$ & $-0.5452$ \\
$3.0$   & $-1.3048$ & $-1.3077$ & $-0.4009$ & $-0.4013$ \\
$2.0$   & $-1.3056$ & $-1.3097$ & $-0.3213$ & $-0.3216$ \\
$1.0$   & $-1.3022$ & $-1.3053$ & $-0.1625$ & $-0.1626$ \\
\bottomrule
\end{tabular}
\end{table}

\paragraph{Cost of the correction.} The factor $(\sigobs/\sigint)^3$ diverges as
$\sigint \to 0$. On the real table, $391$ of $1250$ rows have
$\sigint/\sigphot < 1$ and $34$ have $\sigint = 0$ exactly; computing
\eqref{eq:g1int} without a signal-to-noise gate produces a distribution with
standard deviation $341$ and maximum $1.19\times10^{4}$, against an interquartile
range of only $[-0.25, +0.30]$. \textbf{Equation~\eqref{eq:g1int} must be gated on
$\sigint/\sigphot$ above some threshold} (unity is a reasonable starting choice)
or it will contaminate anything downstream.

% ----------------------------------------------------------------------------
\subsection{The $\max(0,\cdot)$ floor}
\label{sec:floor}
% ----------------------------------------------------------------------------

All three amplitude statistics clamp a negative noise-subtracted variance to
zero. When the intrinsic signal lies at or below the photometric noise the
subtraction goes negative and is truncated, so the quiet end of the sample piles
up at exactly zero rather than scattering about it. This destroys information: the
signed quantity $\sigobs^2 - \sigphot^2$ is an unbiased (if noisy) estimator,
whereas its clamped square root is biased upward. For regression at low
amplitude, retaining the signed variance is preferable. The floor is also what
makes the mask in \eqref{eq:rmsstd} necessary and what makes \eqref{eq:g1int}
undefined on $34$ rows.

% ============================================================================
\section{Robustness to the median normalisation}
% ============================================================================

\subsection{Single-sector rows: an exact affine map}

For a single sector, \eqref{eq:norm} is a global affine transformation
$x \to (f - a)/c$ with $a = 0$, $c = \mathrm{med}(f) > 0$. Behaviour under such a
map partitions the statistics into two classes.

\paragraph{Scale-covariant (amplitudes).} \istd, \irms, \imad{} and $\Delta$ all
carry one power of flux and therefore scale as $c^{-1}$. They are \emph{not}
invariant --- deliberately so. The normalisation is what converts them from
instrumental counts into relative flux, which is the entire reason two clusters
of different apparent brightness can be compared at all.

\paragraph{Scale-invariant (shapes).} $\gp$ and $g_1$ are dimensionless ratios:
%
\begin{align}
  g_1 &\to \frac{\mu_3/c^{3}}{(\sigobs/c)^{3}} = g_1,
  &
  \gp &\to \frac{\Delta/c}{\sigobs/c} = \gp .
  \label{eq:scaleinv}
\end{align}
%
Both are additionally shift-invariant, being built from central quantities.
\textbf{The median normalisation changes neither by any amount.}

This invariance is inherited by the corrected form \eqref{eq:g1int}: because
\eqref{eq:norm} divides the uncertainties by the same $c$, $\sigphot$ and hence
$\sigint$ land in the same relative units as $\mu_3$, so
$\mu_3/\sigint^3$ is scale-invariant as well.

\subsection{Multi-sector rows: a piecewise map}

For a stitched row the normalisation is applied \emph{per sector}, each with its
own divisor $c_k$, and then the sectors are concatenated. This is
\emph{piecewise} affine, not globally affine, so \eqref{eq:scaleinv} does not
apply and every statistic --- shapes included --- depends on the choice.

The choice matters, and per-sector is the correct one. Sectors sit at different
instrumental flux levels (different pixel masks, different CCDs). Stitching
without per-sector normalisation would leave a multi-modal pooled distribution of
$K$ clumps at different levels, with the following effects:

\begin{itemize}
  \item \textbf{Amplitudes are corrupted at first order.} The pooled variance
        acquires the between-clump term $\sum_k w_k (\bar{x}_k - \bar{x})^2$
        directly, so \istd{} would be measuring the instrumental level
        difference rather than the star. \emph{Any} offset inflates it.
  \item \textbf{Shapes are corrupted only at second order.} A multi-modal
        distribution is asymmetric only if the clumps are \emph{unequally
        weighted}. Two equal-length sectors offset from one another produce a
        perfectly symmetric bimodal density: variance badly corrupted, skewness
        still exactly zero.
\end{itemize}

Hence per-sector normalisation is load-bearing for both classes, but it rescues
the amplitudes from a first-order error and the shapes from a second-order one.
Contrary to intuition, the amplitudes are the \emph{more} exposed of the two.

% ============================================================================
\section{Robustness to sector stitching}
% ============================================================================

Once the per-sector normalisation has removed the offsets, a residual and
entirely separate issue remains: a stitched row is a statistical \emph{mixture}
of its sectors, and the statistics survive mixing differently.

\subsection{Pooled moments}

Let sector $k$ contribute $n_k$ points with weight $w_k = n_k/N$, per-sector mean
$\bar{x}_k$, variance $\sigma_k^2$, third moment $\mu_{3,k}$, and offset
$\Delta_k = \bar{x}_k - 1$. Exactly,
%
\begin{align}
  \sigobs^{2} &= \sum_k w_k \left[\sigma_k^{2}
                 + (\bar{x}_k - \bar{x})^{2}\right],
  \label{eq:poolvar}\\
  \mu_3       &= \sum_k w_k \left[\mu_{3,k}
                 + 3\sigma_k^{2}(\bar{x}_k - \bar{x})
                 + (\bar{x}_k - \bar{x})^{3}\right].
  \label{eq:poolmu3}
\end{align}
%
Because per-sector normalisation forces $\mathrm{med}(x_k) = 1$ for every $k$,
the sector means satisfy $\bar{x}_k = 1 + \Delta_k$ with $|\Delta_k| \ll
\sigma_k$ in practice, so the between-sector terms are second order and to
leading order both pooled moments \textbf{mix linearly}:
%
\begin{equation}
  \sigobs^{2} \simeq \sum_k w_k \sigma_k^{2},
  \qquad
  \mu_3 \simeq \sum_k w_k \mu_{3,k},
  \qquad
  \Delta \simeq \sum_k w_k \Delta_k .
  \label{eq:linearmix}
\end{equation}

\subsection{Consequences per statistic}

\paragraph{Amplitudes: well-behaved.} $\istd = \sqrt{\sigobs^2 - \sigphot^2}$ is a
monotone function of a clean linear mixture. A three-sector \istd{} is a
\emph{better-averaged estimate of the same underlying quantity} as a one-sector
\istd{} --- lower epoch-to-epoch scatter, identical meaning. Mixing is benign.

\paragraph{Shapes: reweighted.} Substituting \eqref{eq:linearmix} into
\eqref{eq:g1} and \eqref{eq:gammap} and writing $\mu_{3,k} = g_{1,k}\sigma_k^3$,
$\Delta_k = \gamma_{{\rm P},k}\sigma_k$:
%
\begin{align}
  g_1^{\mathrm{pooled}}
    &= \frac{\sum_k w_k\, g_{1,k}\,\sigma_k^{3}}
            {\left(\sum_k w_k \sigma_k^{2}\right)^{3/2}},
  \label{eq:g1pool}\\[6pt]
  \gp^{\mathrm{pooled}}
    &= \frac{\sum_k w_k\, \gamma_{{\rm P},k}\,\sigma_k}
            {\left(\sum_k w_k \sigma_k^{2}\right)^{1/2}} .
  \label{eq:gppool}
\end{align}
%
Neither is the weighted mean of its per-sector values. The effective weight on
sector $k$ is
%
\begin{equation}
  W_k^{(g_1)} \propto w_k \sigma_k^{3},
  \qquad
  W_k^{(\gp)} \propto w_k \sigma_k^{1},
  \label{eq:effweights}
\end{equation}
%
so the loudest sector dominates --- \emph{cubically} for $g_1$ and only
\emph{linearly} for $\gp$. A stitched skewness is therefore not a better estimate
of the single-sector quantity; it is a \textbf{different quantity}, and how
different depends on the exponent of $\sigma$ in the denominator.

\paragraph{Practical implication.} The cluster table deliberately contains all
seven sector subsets ($r = 1,2,3$) per cluster. Any model trained across
$n_{\rm sec}$ will see the reweighting of \eqref{eq:effweights} as real variation.
Amplitudes are safe; $\gp$ is mildly affected; $g_1$ is the most affected. If
$g_1$ is to be compared across $n_{\rm sec}$, it should be computed \emph{per
sector} and combined explicitly via \eqref{eq:g1pool} rather than evaluated on
the stitched array.

% ============================================================================
\section{Sampling error under the null}
% ============================================================================

Both shape statistics have non-zero scatter on pure noise at finite $N$. For a
Gaussian sample of size $N$ the standard errors are
%
\begin{equation}
  \mathrm{sd}(\gp) = \sqrt{\frac{\pi/2 - 1}{N}} \simeq \frac{0.756}{\sqrt{N}},
  \qquad
  \mathrm{sd}(g_1) = \sqrt{\frac{6}{N}} ,
  \label{eq:nullse}
\end{equation}
%
both confirmed by simulation to three significant figures over
$N \in [500, 4000]$. Dividing by these gives a per-row significance,
%
\begin{equation}
  z_{\gp} = \gp \sqrt{\tfrac{N}{\pi/2-1}},
  \qquad
  z_{g_1} = g_1 \sqrt{\tfrac{N}{6}},
  \label{eq:zscores}
\end{equation}
%
which is likely a better model input than the raw value, since it self-normalises
for the very different $N$ of one-, two-, and three-sector rows.

On the real table (median $N = 1260$), the observed scatter exceeds the null by a
factor of $2.94$ for $\gp$ ($0.0561$ observed versus $0.0191$ expected) and
$4.58$ for $g_1$ ($0.2825$ versus $0.0618$). Significant asymmetry at
$|z| > 3$ is detected in $258/1250$ rows for $\gp$ and $467/1250$ for $g_1$.
\textbf{Real asymmetry is present and $g_1$ is the more sensitive detector},
despite the $n_\sigma = 4$ clipping discussed below.

% ============================================================================
\section{Summary}
% ============================================================================

\begin{table}[htbp]
\centering
\small
\caption{The six statistics along the three axes. ``Median norm.'' refers to a
single global application of \eqref{eq:norm}; for stitched rows the per-sector
application affects all six (\S5.2).}
\label{tab:summary}
\begin{tabular}{llccc}
\toprule
Statistic & Built from & Noise-correctable & Median norm. & Sector mixing \\
\midrule
\istd  & cumulant $\kappa_2$      & exact           & covariant ($c^{-1}$) & benign (linear) \\
\irms  & $\kappa_2$ + quantile    & exact only if $\Delta\!=\!0$ & covariant ($c^{-1}$) & benign (linear) \\
\imad  & quantile                 & approximate     & covariant ($c^{-1}$) & benign (linear) \\
\moff  & $\kappa_1$ + quantile    & no              & covariant ($c^{-1}$) & benign (linear) \\
$\gp$  & $\kappa_1$ + quantile    & no              & \textbf{invariant}   & reweighted $\propto\sigma^{1}$ \\
$g_1$  & cumulant $\kappa_3$      & \textbf{exact}  & \textbf{invariant}   & reweighted $\propto\sigma^{3}$ \\
\bottomrule
\end{tabular}
\end{table}

\subsection{The two skewness statistics are complementary}

$\gp$ and $g_1$ trade off along two \emph{independent} axes, in opposite
directions:

\begin{center}
\begin{tabular}{lcc}
\toprule
 & noise-correctable & survives sector mixing \\
\midrule
$\gp$ & no             & better ($\sigma^1$) \\
$g_1$ & yes, exactly   & worse ($\sigma^3$) \\
\bottomrule
\end{tabular}
\end{center}

Neither dominates. $\gp$ is outlier-robust and survives the pipeline's symmetric
$4\sigma$ clipping intact, but is diluted by noise in a way no denominator can
repair. $g_1$ is exactly de-biasable and empirically the more sensitive detector,
but rests on cubed residuals and is the more distorted by sector mixing. Carrying
both is defensible.

\subsection{Caveat: symmetric sigma clipping}

Step 2 of the pipeline runs \texttt{mad\_clip\_lc} at $n_\sigma = 4$,
\emph{symmetrically} about the sector median. Skewness lives in the tails, so a
symmetric cut actively removes the signal both shape statistics are built to
measure --- flares especially. Values of $\gp$ and $g_1$ computed downstream are
therefore the residual asymmetry of an already partly symmetrised distribution,
and should be read as \emph{lower bounds} on the true asymmetry. If either
statistic proves useful, recomputing at looser $n_\sigma$ (or before the clip) is
the natural follow-up.

\subsection{Empirical status}

On the current table ($1250$ rows, $348$ clusters), neither shape statistic shows
a monotonic relation with cluster age: Spearman $|\rho| \le 0.03$ for $\gp$,
$g_1$, $g_1^{\mathrm{(int)}}$, and both $z$-scores when restricted to the $859$
rows with $\sigint/\sigphot > 1$. (An apparent $\rho = +0.081$ for
$g_1^{\mathrm{(int)}}$ over all rows is an artifact of the divergence described
in \S\ref{sec:g1}; it vanishes under the SNR cut.) Asymmetry is detected at high
significance, but it does not track age.

A plausible physical explanation is that these are \emph{integrated} cluster
light curves --- one flux time series per cluster, summing many stars. Individual
spot and flare signatures, which are precisely the asymmetric features these
statistics detect, average toward symmetry by the central limit theorem as more
stars are blended. This is testable: the scatter in $\gp$ should decrease with
cluster membership count. If it does, skewness is the wrong instrument for
blended photometry irrespective of the clipping; if it does not, the $n_\sigma=4$
clip becomes the prime suspect.

\end{document}
